% s8: from the rooms to all of space (point 5, second half). ASCII only.
\section{From the rooms to all of space}

One road builds the theory in finite boxes and then takes their infinite-volume limit. I start from spheres and stay on spheres: I propose to define the theory on $\mathbb R^4$ as
the $R\to\infty$ end of the continuum theories on $S^3_R\times\mathbb R$. Every ball of fixed size in a sphere of radius
$R$ becomes flat as $R$ grows, providing a candidate geometric limit, not convergence of quantum correlations by itself.
A separate box-versus-sphere identification is unnecessary for this chosen route. In my register, line 29 is folded into line 23,
and two hypotheses become one, H4$'$~\cite{ClosedSet}. Clay asks for no box, only a theory on $\mathbb R^4$ with the
Osterwalder--Schrader properties~\cite{JaffeWitten}.

Inside each room the fine grain must still be removed. The templates are Ba{\l}aban's block-spin renormalization
group~\cite{Balaban} and the construction of Magnen, Rivasseau and S\'en\'eor~\cite{MRS93}.

\begin{claim}{PROVED (CITED)}{The ultraviolet limit in a finite box (partial)}\label{b:box}
Four-dimensional lattice Yang--Mills in a finite box is ultraviolet stable uniformly in the spacing, and the $SU(2)$
Schwinger functions exist with the ultraviolet cutoff removed at fixed infrared cutoff.
\evidence{\cite{Balaban,MRS93}. Neither removes the infrared cutoff or proves a gap; on my geometry they are templates
for the sphere, not results on it.}
\lines{21}
\end{claim}

\begin{claim}{PROVED}{Reflection positivity survives limits}\label{b:rplimit}
$\langle\overline{\Theta F}\,F\rangle\ge0$ for observables $F$ on one side of a reflection is a closed condition, so it
holds in every limit of reflection-positive measures in which these expectations converge.
\evidence{A limit of nonnegative numbers is nonnegative~\cite{GlimmJaffe}; the lattice input is Statement~\ref{b:rp}. For flowed
observables, which are not strictly local, an approximation lemma is needed, which I place inside H4$'$.}
\lines{22, 2}
\end{claim}

Two features are new on the sphere. Nothing mixes time with space on $S^3\times\mathbb R$, so $\mathrm{tr}\,F_{0i}^2$
and $\mathrm{tr}\,F_{ij}^2$ renormalize separately and the limit tunes \emph{two} bare parameters, the coupling $g_0$ and a
speed of light $c_0$. And bare Wilson loops diverge with their perimeter in four dimensions, so I measure observables
smoothed by the gradient flow, perturbatively finite at positive flow time after ordinary renormalization
\cite{Luscher10,LuscherWeisz11}. Flow smears in time as well as space, so a bounded flowed loop is not automatically
a time-slice Hilbert-space operator or an admissible reflection-positive observable.

\begin{claim}{CONJECTURE}{H4$'$: the sphere family (line 23, with line 29 merged)}\label{b:h4}
For every $R$ the ultraviolet limit exists on $S^3_R\times\mathbb R$, with the cutoff of H2 removed, $(g_0,c_0)$ tuned,
and Osterwalder--Schrader regularity bounds. As $R\to\infty$ the flowed correlation functions converge on compact regions
to an $\mathbb R^4$ theory with full Euclidean invariance. This hypothesis also includes a controlled local or
reflection-compatible observable algebra, density of its reconstructed states, and uniform bounds through the limits.
Positive-flow data must be connected to that algebra, not substituted for locality. Given these ingredients and
the remaining OS properties (clustering supplied with H3 below), reconstruction gives a Wightman theory~\cite{OS73,OS75,GlimmJaffe}.
\evidence{I have not constructed this family. Support: the box results~\cite{Balaban,MRS93}; the computed agreement
with $b_0$ (Statement~\ref{b:b0}); perturbative positive-flow finiteness~\cite{Luscher10,LuscherWeisz11}; the
flat-limit derivative is a flowed stress-tensor insertion~\cite{Suzuki13}.}
\lines{23, 29}
\end{claim}

\textbf{The gap survives the limit.} For the gap, the step from rooms to space is one clean lemma. By the spectral
theorem, for $v_A=(A-\langle\Omega,A\Omega\rangle)\Omega$, the autocorrelation
$\langle v_A,e^{-\tau H}v_A\rangle=\int_{[0,\infty)}e^{-\tau E}d\nu_A(E)$ has a positive spectral measure.
An atom at zero is not excluded in advance.

\begin{claim}{PROVED}{The spectral lemma}\label{b:lemma}
If $0\le\langle v_A,e^{-\tau H}v_A\rangle\le C_Ae^{-m\tau}$ for all $\tau>0$ and a set of $v_A$ dense in
$\Omega^\perp$, with finite $C_A$, then
$\mathrm{spec}\,H\cap(0,m)=\emptyset$.
\evidence{Weight of $\nu_A$ on $[0,m)$ would decay more slowly than $e^{-m\tau}$; density does the rest~\cite{ReedSimonIV}.}
\lines{28}
\end{claim}

\begin{claim}{CONDITIONAL}{The gap reaches $\mathbb R^4$, with exponential clustering}\label{b:deduce}
For bounded physical operators $A_R$ with uniformly bounded operator norms, H3 gives the bound with rate
$c\Lambda$ and $C_A\le\sup_R\|A_R\|^2$. If H4$'$ identifies their correlations and supplies dense limiting states,
the bound passes to the limit and the spectral lemma gives $\mathrm{spec}\,H\cap(0,c\Lambda)=\emptyset$.
Locality and Euclidean symmetry then give spatial clustering, with decay rates below the gap and appropriate
observable-dependent constants. A classical sup norm of an arbitrary four-dimensionally flowed observable is not
by itself the required operator-norm or density argument.
\evidence{On H3 and H4$'$. The pass to the limit runs along the sphere family and uses only the sphere-side gap, never
the flat-space result it is proving.}
\lines{28, 20, 23}
\end{claim}

\begin{claim}{CONDITIONAL}{A unique vacuum}\label{b:vacuum}
The $\mathbb R^4$ vacuum is unique: the spectral measure of every truncated state has no atom at zero.
\evidence{It follows from the hypothesis of the lemma, $\nu_A(\{0\})=0$, not from its conclusion; in the reduced model it
holds by Perron--Frobenius~\cite{GlimmJaffe,ReedSimonIV}. So it is conditional on exactly what Statement~\ref{b:deduce} is.}
\lines{27, 28}
\end{claim}

\textbf{An idea for the time--space residue.} Round $S^3\times\mathbb R$ is conformally $\mathbb R^4$ minus a point,
and classical Yang--Mills is conformal; the quantum breaking is the trace anomaly, $\propto b_0\,g^2\,\mathrm{tr}\,F^2$
\cite{CDJ77}, whose size is the $b_0$ I matched in \S6.

My conjecture, an idea rather than an argument, is that the separate tuning of $c_0$ is controlled entirely by this
anomaly, so full Euclidean invariance at $R\to\infty$ follows from the conformal map and the matched $b_0$ (cleanest at
$t=1$; $S^3_t$ is not conformally flat, though every fixed ball still flattens).

If H3 and H4$'$ hold, the limit is a relativistic quantum theory on $\mathbb R^4$ with one vacuum and nothing between it
and $c\Lambda$. That is existence and a gap. The Clay statement also asks what the theory contains.
